\subsection{Existence of representatives}

PROPOSITION 3. Let K be a (not necessarily commutative) locally compact field which is not discrete and whose topology is defined by a discrete valuation v; let A be the ring and m the ideal of v and let us write \(\operatorname{Card}(\mathbf{A}/\mathfrak{m}) = q = p^{f}\) (p prime). Then, there exists a system of representatives S of A/m in A and a uniformizer u for v such that \(0 \in S\) , \(S^{*} = S \cap K^{*}\) is a cyclic subgroup of \(K^{*}\) and \(u^{-1} Su = S\) . Moreover, every element of A may be written uniquely in the form \(\sum_{i=0}^{\infty} s_{i} u^{i}\) , where \(s_{i} \in S\) .

We shall use the following lemma:

LEMMA 1. Let \(x, y\) be two permutable elements of \(A\) such that \(x - y \in \mathfrak{m}^j\) ( \(j \geqslant 1\) ); then \(x^{p^n} - y^{p^n} \in \mathfrak{m}^{j + n}\) for every integer \(n \geqslant 0\) .

By induction on \(n\) , it is reduced to proving the lemma for \(n = 1\) . Then

\[
x ^ {p} - y ^ {p} = (x - y) \left(x ^ {p - 1} + x ^ {p - 2} y + \dots + y ^ {p - 1}\right);
\]

the second factor is a sum of \(p\) terms congruent to one another mod. \(m\) and, as \(A / m\) is of characteristic \(p, p\) . \(1 \in m\) in \(A\) and hence

\[
x ^ {p - 1} + x ^ {p - 2} y + \dots + y ^ {p - 1} \in \mathfrak {m};
\]

whence \(x^{p} - y^{p}\in \mathfrak{m}^{j + 1}\)

We know that the multiplicative group \((\mathbf{A} / \mathfrak{m})^*\) is a cyclic group with \(q - 1\) elements (Algebra, Chapter V, § 11, no. 1, Theorem 1); let \(x\) be a representative in A of a generator of this group; then \(x^{q} - x \in \mathfrak{m}\) , whence, by Lemma 1, \(x^{q^{n + 1}} - x^{q^n} \in \mathfrak{m}^{1 + fn}\) , since \(x^{q}\) and \(x\) are permutable. This proves that \((x^{q^n})_{n \geqslant 0}\) is a Cauchy sequence in A; as A is compact and hence complete, this sequence has a limit \(s\) in A which obviously satisfies \(s \equiv x\) (mod. \(\mathfrak{m}\) ) and \(s^q = s\) . As \(s \neq 0\) , \(s^{q - 1} = 1\) , more precisely \(s\) is a primitive \((q - 1)\) -th root of unity in A. Clearly the set S consisting of 0 and the powers \(s^j\) ( \(0 \leqslant j \leqslant q - 2\) ) is a system of representatives of the classes of A mod. \(\mathfrak{m}\) and is invariant under multiplication in A.

Now let \(a\) be a uniformizer for \(v\) and consider the inner automorphism \(y \mapsto a^{-1}ya\) of \(K\) ; it maps \(A\) to itself, \(m\) to itself and therefore, taking quotients, it defines an automorphism of the field \(A/m\) ; it is known (Algebra, Chapter V, § 11, no. 4, Proposition 5) that such an automorphism is of the form \(z \mapsto z^{p^r}\) , where \(0 \leqslant r \leqslant f - 1\) . Then \(a^{-1}s^j a \equiv s^{jp^r} \pmod{m}\) for \(0 \leqslant j \leqslant q - 2\) ; as \(a \in m\) and \(s \notin m\) , this implies that \(s^{-j}as^{jp^r} \equiv a \pmod{m^2}\) .

Let us write

\[
u = \sum_ {j = 0} ^ {q - 2} s ^ {- j} a s ^ {j p ^ {r}}.
\]

Then \(u \equiv (q - 1)a \equiv -a \pmod{m^2}\) since \(p.1 \in m\) ; we conclude that \(u\) is also a uniformizer for \(v\) ; moreover

\[
s ^ {- 1} u s ^ {p ^ {r}} = u\tag{3}
\]

whence we deduce that \(u^{-1} \mathrm{S}u = \mathrm{S}\) .

Finally, for all \(x \in \mathbf{A}\) there exists a unique sequence \((s_i)\) ( \(i \in \mathbf{N}\) ) such that \(s_i \in \mathbf{S}\) for all \(i\) and \(x \equiv \sum_{i=0}^{n} s_i u^i \pmod{\mathfrak{m}^{n+1}}\) for all \(n \geqslant 0\) : it is immediate by induction on \(n\) , every element \(t\) of \(\mathfrak{m}^{n+1}\) satisfying a relation of the form \(t \equiv t'u^{n+1} \pmod{\mathfrak{m}^{n+2}}\) , where \(t'\) is a uniquely determined element of \(S\) . Then \(x = \sum_{i=0}^{\infty} s_i u^i\) and the family \((s_i)\) satisfying this relation and such that \(s_i \in S\) for all \(i\) is determined uniquely.

